\section{PKI, Chain of Trust, dan Protokol TLS}

\subsection{Peran Diffie-Hellman dalam TLS}

Diffie-Hellman memungkinkan dua pihak menyepakati kunci rahasia melalui kanal tidak aman \cite{stallings2017}. Parameter publik: prima $p$, generator $g$. Keamanan bergantung pada kesulitan logaritma diskrit (diperkenalkan di Bab 6).

\begin{enumerate}
  \item Alice memilih $a$ rahasia, mengirim $A = g^a \bmod p$ ke Bob
  \item Bob memilih $b$ rahasia, mengirim $B = g^b \bmod p$ ke Alice
  \item Alice menghitung $K = B^a \bmod p$
  \item Bob menghitung $K = A^b \bmod p$
  \item Keduanya memperoleh $K$ yang sama
\end{enumerate}

\subsection{TLS (Transport Layer Security)}

TLS mengamankan komunikasi di internet (HTTPS) \cite{stallings2017}. Menggabungkan pertukaran kunci (DH atau ECDH), autentikasi sertifikat, dan enkripsi simetris (AES-GCM, ChaCha20-Poly1305). Handshake: negosiasi cipher suite, verifikasi sertifikat server, pertukaran kunci sesi, lalu enkripsi aplikasi.

\begin{itemize}
  \item \textbf{Negosiasi cipher suite}: Server dan klien menegosiasi algoritma yang didukung.
  \item \textbf{Autentikasi server}: Server membuktikan identitas dengan sertifikat digital.
  \item \textbf{Pertukaran kunci}: Diffie-Hellman atau RSA untuk pertukaran kunci sesi.
  \item \textbf{Enkripsi data}: AES-GCM atau ChaCha20-Poly1305 untuk data.
  \item \textbf{Integritas}: AEAD (GCM) atau HMAC untuk integritas.
\end{itemize}

\subsection{Chain of Trust dan Sertifikat}

TLS membutuhkan \textbf{otentikasi server} melalui sertifikat X.509 \cite{stallings2017}. \textbf{Chain of trust}: sertifikat server ditandatangani CA menengah; CA menengah ditandatangani CA root. Klien memvalidasi rantai sampai CA root dalam daftar tepercaya (trust store). Tanpa validasi sertifikat, koneksi rentan man-in-the-middle.

TLS 1.2 dan TLS 1.3 adalah versi saat ini. Forward secrecy: kunci sesi ephemeral sehingga kebocoran kunci privat tidak membeberkan sesi lampau \cite{stallings2017}.
